\documentclass[12pt,a4paper,twoside]{article}
\usepackage{hyperref}

% --- Sprache, Kodierung & Typografie ---
\usepackage[utf8]{inputenc}       % UTF-8 Kodierung
\usepackage[T1]{fontenc}          % Korrekte Silbentrennung & Umlaute
\usepackage[french]{babel}        % Spracheinstellungen
\usepackage{microtype}            % Bessere Typografie (Zeilenabstände, Umbrüche)

% --- Schrift ---
\usepackage{mathpazo}             % Palatino für Text & Mathe (sehr lesefreundlich)

% --- Seitenlayout ---
\usepackage[a4paper,margin=2.8cm]{geometry}
\setlength{\headheight}{15pt}

% --- Mathematik & Symbole ---
\usepackage{amsmath, amssymb, amsthm, mathtools}
\usepackage{bm}                   % Fette mathematische Symbole
\usepackage{latexsym}

% --- Code ---
\usepackage{listings}
\usepackage{xcolor}

\definecolor{codegreen}{rgb}{0,0.6,0}
\definecolor{codegray}{rgb}{0.5,0.5,0.5}
\definecolor{codepurple}{rgb}{0.58,0,0.82}
\definecolor{backcolour}{rgb}{0.95,0.95,0.92}

\lstdefinestyle{mystyle}{
    backgroundcolor=\color{backcolour},   
    commentstyle=\color{codegreen},
    keywordstyle=\color{blue},
    numberstyle=\tiny\color{codegray},
    stringstyle=\color{codepurple},
    basicstyle=\ttfamily\footnotesize,
    breakatwhitespace=false,         
    breaklines=true,                 
    captionpos=b,                    
    keepspaces=true,                 
    numbers=left,                    
    numbersep=5pt,                  
    showspaces=false,                
    showstringspaces=false,
    showtabs=false,                  
    tabsize=2
}

\lstset{style=mystyle}

% --- Grafiken & Zitate ---
\usepackage[pdftex]{graphicx}
\usepackage{url}
\usepackage{cite}
\usepackage{float}
\bibliographystyle{plain}         % oder "alpha" / "abbrv" / "apalike"

%--- Beschreibungen ---
\usepackage{caption}
\DeclareCaptionFormat{custom}{\textbf{#1#2}#3}  %Fetter name "Figure ..."
\captionsetup{
    format=custom,
    labelformat=original,
    labelsep=period,
    justification=centering %Damit alle Zeilen in der Mitte sind
}

% --- Zeilenabstand ---
\usepackage{setspace}
\onehalfspacing                   % 1.5-facher Zeilenabstand (Standard bei Masterarbeiten)

% --- Kopf- und Fußzeilen ---
\usepackage{fancyhdr}
\pagestyle{fancy}
\fancyhf{} % clear header/footer

% --- Headers ---
\fancyhead[LE]{\thepage}               % Seite links (Even pages)
\fancyhead[RO]{\thepage}               % Seite rechts (Odd pages)
\fancyhead[RE,LO]{\nouppercase{\leftmark}} % Kapitel-/Abschnittstitel innen

% --- Linien ---
\renewcommand{\headrulewidth}{0.4pt}         % Linie unter Kopfzeile
\renewcommand{\footrulewidth}{0pt}           % keine Linie unten

% --- Theorem-Umgebungen ---
\newtheorem{Theorem}{Theorem}[section]
\newtheorem{Definition}[Theorem]{Definition}%[section]
\newtheorem{Lemma}[Theorem]{Lemma}
\numberwithin{equation}{section}

% --- Abkürzungen ---
\newcommand{\f}{\varphi}
\newcommand{\Poly}{\mathbb{P}}
\newcommand{\N}{\mathbb{N}}

% no indentation
\setlength{\parindent}{0em}

\begin{document}
  % no page numbers 
  \pagestyle{empty}

  % title page
\begin{titlepage}
    \begin{center}

    % title
    \vspace*{1cm}
    {\Large \textbf{Projet Python : Animation mathématique}}\\[1.5cm]

    {\large  LU3MA201 - Projet / Travail d'étude et de recherche}\\
    {\large  Licence 3 Mathématiques}\\
    {\large  Sorbonne Université}\\[2cm]

    {\Large Rapport de projet}\\[2cm]

    {\large Solveig GIRARDIN, Walid KHALFALLAH,\\ Anis LAPLAGNE et Yuyang ZHOU}\\[1cm]
    {\large 30/05/2025}\\[3cm]

 \end{center}
\end{titlepage}

\clearpage
% table of contents
\tableofcontents
\newpage

% fancy headers
\pagestyle{fancy}

\section{Introduction}
Le Youtubeur Grant Sanderson, qui anime la chaîne de vulgarisation
\href{https://www.youtube.com/3Blue1Brown}{3Blue1Brown}, a créé une bibliothèque
Python, appelé \textit{manim} (pour "animations mathématiques"), pour réaliser
les nombreuses animations utilisées dans ses vidéos. Le but du projet est de
l'utiliser pour créer des animations illustrant des concepts du cours de
topologie de L3.

\clearpage
\section{Fonctionnement de \textit{manim} sur un  exemple simple}
Pour expliquer le fonctionnement de la bibliothèque \textit{manim}, nous allons
prendre l'exemple de la définition d'un ensemble ouvert. 
\begin{Definition}
Soit $(X,d)$ un espace métrique. Une partie O de $X$ est dite
\underline{ouverte} si tout point de O est le centre d'une boule ouverte incluse
dans O. Autrement dit :
$$\forall x \in O \;\exists \varepsilon > 0 \text{ t.q. } B(x,\varepsilon) \subset O.$$
\end{Definition}

Pour animer cette définition d'une façon simple, l'idée est de représenter
l'ensemble $O$ par un cercle (bleu) sans bord et d'y placer un point $x$
arbitraire. Depuis ce point, va alors grossir un second cercle, la boule ouverte
$B(x, \varepsilon)$, jusqu'à ce que son rayon, $\varepsilon$, atteint la limite
supérieure des rayons des boules centrées en $x$ contenues dans $O$.\\

Pour commencer, il faut importer tous les éléments de la bibliothèque
\textit{manim}. Ensuite, chaque séquence d'animation que nous voulons produire
doit être définie dans un bloc de code commençant par le mot-clé \verb+class+
suivi d'un nom choisi pour cette séquence et du mot-clé \verb+Scene+ (ou dans
certains cas \verb+ZoomedScene+) entre parenthèses. Dans notre cas, cela nous
donne le code suivant :\\

\begin{lstlisting}[language=Python]
from manim import *
class DefOuvert(Scene):
\end{lstlisting}

Maintenant nous devons instancier les objets géométriques (cercles, points, segments,\ldots) que nous voulons animer dans un bloc de code indenté en-dessous des mots-clés \verb+def constructor(self)+. Ici, nous aurons besoin de deux cercles, d'un point $x$ et d'un segment qui illustre le rayon $\varepsilon$ de la boule.\\

Plus précisément, nous allons définir une variable \verb|pos_ensemble| avec laquelle nous allons déterminer la position du centre de notre ensemble $O$ (Ici nous avons choisi le milieu de l'écran donc les coordonnées 0,0,0)) et une deuxième variable \verb|rayon_ensemble| pour la valeur de son rayon.
Nous obtenons alors :\\

\begin{lstlisting}[language=Python]
from manim import *
from math import sqrt # Pour le calcul de la valeur d'epsilon

class DefOuvert(Scene):
    def constructor(self):
        # Ensemble O
        pos_ensemble = np.array([0,0,0])
        rayon_ensemble = 1.5
        # On construit un cercle bleu de rayon "rayon_ensemble" et de facteur d'opacite 0.25
        # Puis on positionne son centre au point "pos_ensemble"
        ensemble = Circle(radius = rayon_ensemble, color = BLUE, fill_opacity = 0.25).move_to(pos_ensemble)
        ensemble.set_stroke(width = 0)  # Enleve la bordure du cercle

        # Point x
        position = np.array([-0.15,-0.2,0])
        point = Dot(color = YELLOW).rayon(0.5).move_to(position)

        # Boule
        epsilon = rayon_ensemble - sqrt(((pos_ensemble-pos_p1)**2).sum())
        boule = Circle(radius=epsilon, color=GREEN, fill_opacity = 0.5).move_to(position)
        boule.set_stroke(width = 0) # Enleve la bordure du cercle

        # Rayon epsilon
        rayon = Line(point.get_center(), boule.point_at_angle(7*PI/8)).set_color(RED).set_z_index(0.5)
        
\end{lstlisting}

Maintenant, nous avons tous les objets mais ils n'aparaissent pas encore. Pour
cela, nous devons faire appel à la méthode \verb+self.play()+ (héritée de
\verb+Scene+) avec toutes les animations que nous voulons faire en même temps.
Il a un certain nombre d'animations prédéfinies que nous pouvons simplement
appeler en leur passant l'objet que nous voulons animer de cette manière-là en
argument. Concrètement, si nous voulons faire apparaître notre ensemble $O$ et
point $x$ en même temps à l'aide de l'animation \verb+FadeIn()+ (ils vont
simplement graduellement apparaître sur l'écran), il faut ajouter la ligne
suivante à la fin de notre constructeur :\\

\begin{lstlisting}[language=Python]
        self.play(FadeIn(ensemble), FadeIn(point))
\end{lstlisting}

Avant de faire apparaître la boule et son rayon, nous pouvons faire une pause
dans l'animation avec la méthode \verb|self.wait()|. Sa valeur par défaut est
une seconde, mais si nous souhaitons attendre plus longtemps, il suffit de
l'appeler avec un temps en secondes. \\

Pour la boule et le rayon, il est intéressant de les faire grandir à partir du
point pour bien illustrer le fait qu'ils dépendent du point $x$ choisit et de sa
position à l'intérieur de l'ensemble. L'animation à choisir pour obtenir cet
effet est \verb|GrowFromCenter()| pour la boule et \verb|GrowFromPoint()| pour
le rayon.\\

Le code final que nous obtenons est alors le suivant. \\

\begin{lstlisting}[language=Python]
from manim import *
from math import sqrt # Pour le calcul de la valeur d'epsilon

class DefOuvert(Scene):
    def constructor(self):
        # Ensemble O
        pos_ensemble = np.array([0,-1,0])
        rayon_ensemble = 1.5
        ensemble = Circle(radius = rayon_ensemble, color = BLUE, fill_opacity = 0.25).move_to(pos_ensemble)
        ensemble.set_stroke(width = 0)  # Enleve la bordure du cercle

        # Point x
        position = np.array([-0.15,-0.2,0])
        point = Dot(color = YELLOW).rayon(0.5).move_to(position)

        # Boule
        epsilon = rayon_ensemble - sqrt(((pos_ensemble-pos_p1)**2).sum())
        boule = Circle(radius=epsilon, color=GREEN, fill_opacity = 0.5).move_to(position)
        boule.set_stroke(width = 0) # Enleve la bordure du cercle

        # Rayon epsilon
        rayon = Line(point.get_center(), boule.point_at_angle(7*PI/8)).set_color(RED)

        # Animations
        self.play(FadeIn(ensemble), FadeIn(point))
        self.wait()
        self.play(GrowFromCenter(boule), GrowFromPoint(rayon, rayon.get_start()))
        
\end{lstlisting}

\end{document}